\section{Diffusion Models}
\label{sec:diffusion}

A GAN turns random numbers into an image in a single pass of its generator and
learns this mapping through an adversarial game. A diffusion model reaches the
same goal on a very different path. It first destroys images on purpose: over
many small steps, it mixes a training image with Gaussian noise until nothing
but noise is left. This destruction needs no learning at all. A neural network
then learns to undo one small step at a time. To generate a new image, one
starts from pure noise and applies the learned denoising step again and again
until a clean image appears. Each individual step is an easy regression
problem with a known target, so training is as stable as ordinary supervised
learning. This section derives the forward noising process, the noise
prediction objective, and the reverse sampling procedure, then shows how
Stable Diffusion runs the process in a compressed latent space conditioned on
text, and how low-rank adaptation (LoRA) fine-tunes such a large model on a
small image set.

\subsection{Destroying and Recreating Images}

Figure~\ref{fig:diff-chain} shows the two directions. The \emph{forward} or
\emph{diffusion process} maps the input image $x_0$ through a chain of latent
variables $x_1,x_2,\dots,x_T$. Each step adds a little noise, and after
$T$ steps (typically $T=1000$) the image $x_T$ is indistinguishable from a
sample of a standard normal distribution. In the language of the latent
variable models of Section~\ref{sec:gan}, this chain acts as an
\emph{encoder}, but a fixed one without parameters. The \emph{reverse process}
is the decoder. It passes the data back through the same latent variables and
removes noise at each step. Only the reverse process is learned.

\begin{figure}[ht]
  \centering
  \begin{tikzpicture}[
    img/.style={draw, minimum size=1.3cm, inner sep=0pt},
    fwd/.style={-{Latex[length=2mm]}, thick},
    rev/.style={-{Latex[length=2mm]}, thick, orange!80!black},
    lbl/.style={font=\footnotesize, align=center}
  ]
    \pgfmathsetseed{7}
    \foreach \i/\x/\n in {0/0/0, 1/2.3/25, 2/4.6/70, 3/8.4/260} {
      \node[img] (x\i) at (\x,0) {};
      \ifnum\i<3
        \fill[blue!45] (\x,0.05) circle (0.33);
        \fill[blue!45] (\x-0.4,-0.55) rectangle (\x+0.4,-0.3);
      \fi
      \ifnum\n>0
        \foreach \k in {1,...,\n} {
          \fill[black!70] ($(\x,0)+(rnd*1.2-0.6, rnd*1.2-0.6)$) circle (0.025);
        }
      \fi
    }
    \node at (6.5,0) {$\cdots$};
    \node[lbl, below=0.1cm of x0] {$x_0$\\ clean image};
    \node[lbl, below=0.1cm of x1] {$x_1$};
    \node[lbl, below=0.1cm of x2] {$x_2$};
    \node[lbl, below=0.1cm of x3] {$x_T$\\ $\approx\mathcal N(0,I)$};
    \draw[fwd] ($(x0.north)+(0.2,0.15)$) to[bend left=30]
      node[lbl, above] {$q(x_1\mid x_0)$} ($(x1.north)+(-0.2,0.15)$);
    \draw[fwd] ($(x1.north)+(0.2,0.15)$) to[bend left=30]
      node[lbl, above] {$q(x_2\mid x_1)$} ($(x2.north)+(-0.2,0.15)$);
    \draw[fwd] ($(x2.north)+(0.2,0.15)$) to[bend left=20]
      node[lbl, above] {fixed noising} ($(x3.north)+(-0.2,0.15)$);
    \draw[rev] ($(x3.south)+(-0.2,-0.75)$) to[bend left=20]
      node[lbl, below] {learned denoising} ($(x2.south)+(0.2,-0.75)$);
    \draw[rev] ($(x2.south)+(-0.2,-0.75)$) to[bend left=30]
      node[lbl, below] {$p_\theta(x_1\mid x_2)$} ($(x1.south)+(0.2,-0.75)$);
    \draw[rev] ($(x1.south)+(-0.2,-0.75)$) to[bend left=30]
      node[lbl, below] {$p_\theta(x_0\mid x_1)$} ($(x0.south)+(0.2,-0.75)$);
  \end{tikzpicture}
  \caption{Diffusion model. The forward process (black, top) adds a little
  Gaussian noise per step until the image is pure noise $x_T$. The reverse
  process (orange, bottom) is a learned network that removes noise step by
  step and so turns a noise sample into an image.}
  \label{fig:diff-chain}
\end{figure}

\noindent
Why go through a thousand small steps instead of mapping noise to an image in
one pass, as a GAN does? A single jump from noise to a realistic image is a
hard, highly ambiguous mapping. A small denoising step, in contrast, only has
to make a slightly noisy image a little cleaner. Because every step adds only
a small amount of noise, the reverse of a step is well approximated by a
Gaussian, and a network can predict it with a plain squared-error loss. There
is no discriminator, no minimax game, and no mode collapse: every training
image is noised and must be reconstructed, so the model is pushed to cover the
whole data distribution. Diffusion models belong to the probabilistic
generative models of Figure~\ref{fig:gen-taxonomy}. Their training objective
is a simplified bound on the likelihood, so, like a VAE, they learn an
approximation of the data distribution. Table~\ref{tab:diff-vs-gan-vae}
compares the three families. The price for stable training is slow sampling:
generating one image requires many network evaluations instead of one.

\subsection{The Forward Noising Process}

Let $\mathcal D=\{x_0^{(1)},\dots,x_0^{(N)}\}$ be the training images. The
forward process is controlled by a \emph{noise schedule}, a fixed sequence of
small variances
\begin{equation}
  0<\beta_1<\dots<\beta_T<1,\qquad
  \alpha_t = 1-\beta_t,\qquad
  \bar\alpha_t = \prod_{s=1}^{t}\alpha_s .
  \label{eq:diff-schedule}
\end{equation}
\noindent
A common choice is a linear schedule from $\beta_1=10^{-4}$ to
$\beta_T=0.02$ with $T=1000$. One forward step shrinks the current image a
little and adds fresh Gaussian noise of variance $\beta_t$:
\begin{equation}
  q(x_t\mid x_{t-1}) = \mathcal N\!\left(x_t;\ \sqrt{\alpha_t}\,x_{t-1},\
  \beta_t I\right),
  \qquad\text{i.e.}\qquad
  x_t = \sqrt{1-\beta_t}\,x_{t-1} + \sqrt{\beta_t}\,\varepsilon_t .
  \label{eq:diff-step}
\end{equation}
\noindent
The shrinking factor $\sqrt{1-\beta_t}$ is chosen so that the variance stays
constant. If $x_{t-1}$ has unit variance per pixel, then $x_t$ has variance
$(1-\beta_t)+\beta_t=1$. Without shrinking, the variance would grow with every
step, and the end point would not be a standard normal distribution.

\begin{table}[ht]
  \centering
  \begin{tabular}{@{}p{2.6cm}p{3.6cm}p{3.6cm}p{3.9cm}@{}}
    \toprule
    & GAN & VAE & Diffusion model \\
    \midrule
    Latent space & low-dim.\ $z$ & low-dim.\ $z$ & same size as image,
      $T$ copies \\
    Encoder & none & learned & fixed noising, no parameters \\
    Training & adversarial minimax & likelihood bound & likelihood bound,
      reduced to MSE on noise \\
    Stability & fragile, mode collapse & stable & stable \\
    Density & no & approximate & approximate \\
    Sample quality & sharp & often blurry & sharp, state of the art \\
    Sampling cost & one pass & one pass & many passes ($T$ or ${\approx}50$) \\
    \bottomrule
  \end{tabular}
  \caption{Three families of generative models. Diffusion models combine the
  stable likelihood-based training of VAEs with the sample quality of GANs,
  but pay with many network evaluations per sample.}
  \label{tab:diff-vs-gan-vae}
\end{table}

\begin{figure}[ht]
  \centering
  \begin{tikzpicture}[x=10cm, y=2.6cm,
    lbl/.style={font=\footnotesize, align=left}]
    \draw[-{Latex[length=2mm]}] (0,0) -- (1.06,0) node[right, font=\small] {$t$};
    \draw[-{Latex[length=2mm]}] (0,0) -- (0,1.12);
    \foreach \u/\t in {0/0, 0.25/250, 0.5/500, 0.75/750, 1/1000} {
      \draw (\u,0) -- (\u,-0.03) node[below, font=\footnotesize] {\t};
    }
    \foreach \v in {0, 0.5, 1} {
      \draw (0,\v) -- (-0.01,\v) node[left, font=\footnotesize] {\v};
    }
    % linear schedule: ln abar(u) ~ -(0.1 u + 9.95 u^2), u = t/T
    \draw[thick, blue!70!black, domain=0:1, samples=80]
      plot (\x, {exp(-(0.1*\x + 9.95*\x*\x)/2)});
    \draw[thick, blue!70!black, dashed, domain=0:1, samples=80]
      plot (\x, {sqrt(1 - exp(-(0.1*\x + 9.95*\x*\x)))});
    % cosine schedule: abar(u) = cos^2((u+s)/(1+s)*90deg) / cos^2(s/(1+s)*90deg)
    \draw[thick, orange!85!black, domain=0:1, samples=80]
      plot (\x, {cos((\x+0.008)/1.008*90)/cos(0.008/1.008*90)});
    \node[lbl, blue!70!black] at (0.27,0.40) {$\sqrt{\bar\alpha_t}$, linear};
    \node[lbl, blue!70!black] at (0.85,0.88) {$\sqrt{1-\bar\alpha_t}$, linear};
    \node[lbl, orange!85!black] at (0.73,0.62) {$\sqrt{\bar\alpha_t}$, cosine};
  \end{tikzpicture}
  \caption{Signal coefficient $\sqrt{\bar\alpha_t}$ (solid) and noise
  coefficient $\sqrt{1-\bar\alpha_t}$ (dashed) for the linear schedule with
  $T=1000$, and the signal coefficient of the cosine schedule (orange). The
  linear schedule destroys the image early and wastes many late steps on
  almost pure noise; the cosine schedule spreads the destruction more
  evenly.}
  \label{fig:diff-schedule}
\end{figure}

\paragraph{Closed form for any time step.}
Running a thousand steps to produce one training example would be wasteful.
Fortunately, a sum of independent Gaussians is again Gaussian, so the chain
collapses. Two steps give
\begin{equation*}
  x_t = \sqrt{\alpha_t}\bigl(\sqrt{\alpha_{t-1}}\,x_{t-2}
      + \sqrt{\beta_{t-1}}\,\varepsilon_{t-1}\bigr) + \sqrt{\beta_t}\,\varepsilon_t
      = \sqrt{\alpha_t\alpha_{t-1}}\,x_{t-2}
      + \sqrt{\alpha_t\beta_{t-1}}\,\varepsilon_{t-1}
      + \sqrt{\beta_t}\,\varepsilon_t .
\end{equation*}
\noindent
The two noise terms are independent, so their sum is one Gaussian with
variance $\alpha_t(1-\alpha_{t-1})+(1-\alpha_t)=1-\alpha_t\alpha_{t-1}$. The
two steps thus act like a single step with signal factor
$\sqrt{\alpha_t\alpha_{t-1}}$. By induction over all steps back to $x_0$,
\begin{equation}
  x_t = \sqrt{\bar\alpha_t}\,x_0 + \sqrt{1-\bar\alpha_t}\,\varepsilon,
  \qquad \varepsilon\sim\mathcal N(0,I),
  \qquad\text{i.e.}\qquad
  q(x_t\mid x_0)=\mathcal N\!\left(x_t;\ \sqrt{\bar\alpha_t}\,x_0,\
  (1-\bar\alpha_t)I\right).
  \label{eq:diff-marginal}
\end{equation}
\noindent
This is the key result of the forward process: any noise level $x_t$ can be
sampled directly from $x_0$ with one Gaussian draw. The tensor $\varepsilon$
has the same shape as the image, with one independent Gaussian value per pixel
and channel. It is a \emph{concrete sampled noise tensor}, not an abstract
distribution. The two coefficients split the image into a signal part and a
noise part whose squares always sum to one. The \emph{signal-to-noise ratio}
$\bar\alpha_t/(1-\bar\alpha_t)$ falls monotonically from very large at $t=1$
to almost zero at $t=T$.

\paragraph{Worked example.}
Take one pixel with value $x_0=0.8$ at a time step with $\bar\alpha_t=0.64$,
and suppose the sampled noise is $\varepsilon=-0.5$. Then
$x_t=\sqrt{0.64}\cdot0.8+\sqrt{0.36}\cdot(-0.5)=0.64-0.30=0.34$. For the
linear schedule, $\ln\bar\alpha_T=\sum_t\ln(1-\beta_t)\approx-\sum_t\beta_t$.
The sum of the $1000$ linearly spaced values is
$1000\cdot(10^{-4}+0.02)/2\approx10.05$, so $\bar\alpha_T\approx e^{-10}
\approx4.5\cdot10^{-5}$ and the image coefficient $\sqrt{\bar\alpha_T}$ is
below $0.007$. At the end of the chain almost nothing of $x_0$ survives, and
$q(x_T\mid x_0)\approx\mathcal N(0,I)$ regardless of the input image. This is
exactly the distribution from which generation will start.

\paragraph{Choice of schedule.}
Figure~\ref{fig:diff-schedule} shows that the linear schedule removes most of
the image within the first half of the chain. The \emph{cosine schedule}
defines $\bar\alpha_t$ directly,
$\bar\alpha_t = f(t)/f(0)$ with $f(t)=\cos^2\!\bigl(\tfrac{t/T+s}{1+s}
\cdot\tfrac{\pi}{2}\bigr)$ and a small offset $s=0.008$. It changes slowly near
$t=0$ and $t=T$ and gives better image quality at both ends. More recent
schedules (sigmoid, EDM) are tuned for fast sampling. Whatever the schedule,
it is fixed before training and is not learned.

\subsection{Training a Network to Predict the Noise}

Solving Equation~\ref{eq:diff-marginal} for the noise gives the training
target:
\begin{equation}
  \varepsilon = \frac{x_t-\sqrt{\bar\alpha_t}\,x_0}{\sqrt{1-\bar\alpha_t}} .
  \label{eq:diff-eps-target}
\end{equation}
\noindent
This is the normalized Gaussian noise component that corrupted $x_0$ into
$x_t$. In the worked example, $(0.34-0.64)/0.6=-0.5$ recovers the sampled
value. During training this ground truth is known exactly, because we sampled
$\varepsilon$ ourselves when constructing $x_t$. A neural network
$\varepsilon_\theta(x_t,t)$ receives the noisy image and the time step and
predicts the noise; its loss is the mean squared error
\begin{equation}
  \mathcal L(\theta) = \mathbb E_{x_0,\,t,\,\varepsilon}\Bigl[
  \bigl\|\varepsilon-\varepsilon_\theta\bigl(\sqrt{\bar\alpha_t}\,x_0
  +\sqrt{1-\bar\alpha_t}\,\varepsilon,\ t\bigr)\bigr\|^2\Bigr].
  \label{eq:diff-loss}
\end{equation}
\noindent
The expectation runs over training images, uniformly drawn time steps
$t\in\{1,\dots,T\}$, and noise samples. The full derivation starts from a
bound on the likelihood, as in Equation~\ref{eq:gen-nll}, and contains a
different weight for each $t$. Dropping these weights gives this
\emph{simplified loss}, which works better in practice. One network with
shared parameters $\theta$ serves all noise levels, so it must be
\emph{time-conditioned}. Figure~\ref{fig:diff-train} shows one training step.

\begin{figure}[ht]
  \centering
  \begin{tikzpicture}[
    box/.style={draw, rounded corners=2pt, minimum height=0.9cm,
      align=center, font=\small},
    link/.style={-{Latex[length=2mm]}, thick},
    lbl/.style={font=\footnotesize, align=center}
  ]
    \node[box] (x0) at (0,0.8) {image $x_0$};
    \node[box] (eps) at (0,-0.8) {noise $\varepsilon$};
    \node[box] (t) at (4.8,2.0) {step $t$};
    \node[box, fill=gray!15] (mix) at (3.2,0)
      {$\sqrt{\bar\alpha_t}\,x_0$\\ $+\sqrt{1-\bar\alpha_t}\,\varepsilon$};
    \node[box, fill=orange!25, minimum width=2.2cm] (net) at (6.6,0)
      {U-Net\\ $\varepsilon_\theta(x_t,t)$};
    \node[box, fill=gray!15] (loss) at (10.2,0)
      {$\|\varepsilon-\hat\varepsilon\|^2$};
    \draw[link] (x0) -- (mix);
    \draw[link] (eps) -- (mix);
    \draw[link] (mix) -- node[lbl, above] {$x_t$} (net);
    \draw[link] (net) -- node[lbl, above] {$\hat\varepsilon$} (loss);
    \draw[link] (t.west) -| (mix.north);
    \draw[link] (t.east) -| (net.north);
    \draw[link] (eps.south) -- ++(0,-0.6) -| node[lbl, pos=0.25, below]
      {known target} (loss.south);
  \end{tikzpicture}
  \caption{One training step. A clean image, a noise tensor, and a time step
  are drawn; the mixture $x_t$ goes into the network, and its noise estimate
  $\hat\varepsilon$ is compared with the noise that was actually added.}
  \label{fig:diff-train}
\end{figure}

\noindent
Training thus repeats a short loop:
\begin{enumerate}
  \item draw a training image $x_0\sim\mathcal D$;
  \item draw a time step $t\sim\operatorname{Uniform}\{1,\dots,T\}$;
  \item draw noise $\varepsilon\sim\mathcal N(0,I)$;
  \item form $x_t=\sqrt{\bar\alpha_t}\,x_0+\sqrt{1-\bar\alpha_t}\,\varepsilon$;
  \item compute $\ell=\|\varepsilon-\varepsilon_\theta(x_t,t)\|^2$ and take a
    gradient step on $\theta$.
\end{enumerate}
\noindent
The chain is never run step by step during training. Every example is an
independent regression problem with input $(x_t,t)$ and target
$\varepsilon$.

\paragraph{What the network actually learns.}
The network cannot know the exact noise that was drawn. Many clean images
$x_0$, each combined with a suitable noise tensor, lead to the same noisy image
$x_t$, and at high noise levels a blurry blob is compatible with a cat, a dog,
or a car. For squared error, the best possible prediction is the conditional
mean $\varepsilon_\theta(x_t,t)=\mathbb E[\varepsilon\mid x_t]$, averaged
over all clean images that could have produced $x_t$ under the training
distribution. The network therefore does not memorize an inverse of the
corruption. It learns the statistically best \emph{denoising direction}: the
direction in which $x_t$ must move to become more like a natural image. In
other words, it learns the structure of natural images. This has a precise
counterpart. For the Gaussian of Equation~\ref{eq:diff-marginal},
\begin{equation}
  \nabla_{x_t}\log q(x_t\mid x_0)
  = -\frac{x_t-\sqrt{\bar\alpha_t}\,x_0}{1-\bar\alpha_t}
  = -\frac{\varepsilon}{\sqrt{1-\bar\alpha_t}},
  \label{eq:diff-score}
\end{equation}
\noindent
so predicting the noise is, up to a scale factor, the same as learning the
\emph{score}, the gradient of the log density at each noise level. The score
points uphill in probability, toward regions where natural images lie. This
links diffusion models to score-based generative models. Because $x_t$ and
$\hat\varepsilon$ determine each other, a noise prediction can also be turned
into an estimate of the clean image,
\begin{equation}
  \hat x_0 = \frac{x_t-\sqrt{1-\bar\alpha_t}\,\varepsilon_\theta(x_t,t)}
  {\sqrt{\bar\alpha_t}} ,
  \label{eq:diff-x0-hat}
\end{equation}
\noindent
which at large $t$ is a blurry average of plausible images and at small $t$
is sharp.

\paragraph{The U-Net backbone.}
The noise prediction has the same size as the input, and it must be accurate
both in fine detail and in global layout. The U-Net of
Figure~\ref{fig:unet} meets both requirements, and it is the standard
denoising network. The diffusion U-Net modifies the segmentation U-Net in a
few places. Its encoder consists of residual blocks with spatial
downsampling, for example from $64\times64$ over $32\times32$ and
$16\times16$ down to $8\times8$. At the low resolutions ($16\times16$,
$8\times8$), self-attention layers (Section~\ref{sec:transformers}) capture
global context at moderate cost, since the number of attention pairs grows
with the square of the number of positions. The decoder upsamples and
receives skip connections from the encoder. The time step $t$ is encoded with
sinusoidal features, as in the positional encoding of
Equation~\ref{eq:transformer-position}, passed through a small MLP, and added
to the feature maps of every residual block. The same weights therefore
behave differently at different noise levels: at large $t$ they sketch
coarse layout, at small $t$ they refine texture.

\subsection{Generating Images by Reverse Denoising}

To generate a new image, start from pure Gaussian noise, $x_T\sim\mathcal
N(0,I)$, and apply learned reverse steps
$x_T\to x_{T-1}\to\dots\to x_0$. The true reverse distribution
$q(x_{t-1}\mid x_t)$ is intractable, because it depends on the unknown data
distribution. Since each forward step adds only little noise, it is
approximated by a Gaussian $p_\theta(x_{t-1}\mid x_t)=\mathcal
N(x_{t-1};\mu_\theta(x_t,t),\sigma_t^2I)$, whose mean is computed from the
noise prediction:
\begin{equation}
  x_{t-1} = \frac{1}{\sqrt{\alpha_t}}\left(x_t
  - \frac{\beta_t}{\sqrt{1-\bar\alpha_t}}\,\varepsilon_\theta(x_t,t)\right)
  + \sigma_t z,
  \qquad z\sim\mathcal N(0,I).
  \label{eq:diff-reverse}
\end{equation}
\noindent
Read from the inside out, each step does four things: look at the current
noisy image $x_t$, ask the network to estimate the noise component, subtract
the fraction $\beta_t/\sqrt{1-\bar\alpha_t}$ of that noise, and undo the
shrinking of the forward step with the factor $1/\sqrt{\alpha_t}$. Only a
part of the predicted noise is removed per step, because the network's
estimate is unreliable at high noise; the next step looks again with a better
view. The added term $\sigma_t z$ with fresh noise, usually
$\sigma_t^2=\beta_t$, keeps the process stochastic. It corresponds to
sampling from $p_\theta$ instead of taking its mean, and it lets different
runs from the same start produce different details. At the last step no noise
is added ($z=0$ for $t=1$), and $x_0$ is returned.

\paragraph{Sanity check at $t=1$.}
At the first step $\bar\alpha_1=\alpha_1=1-\beta_1$, so the coefficient in
Equation~\ref{eq:diff-reverse} becomes
$\beta_1/\sqrt{\beta_1}=\sqrt{\beta_1}$, and without the noise term
$x_0=(x_1-\sqrt{\beta_1}\,\varepsilon_\theta)/\sqrt{\alpha_1}$. This is the
exact inverse of the forward step $x_1=\sqrt{\alpha_1}\,x_0+\sqrt{\beta_1}\,
\varepsilon$ of Equation~\ref{eq:diff-step}. If the network predicted the
noise perfectly, the last step would restore the image exactly. For larger
$t$ the formula removes only the share of the total noise that belongs to one
step.

\paragraph{Sampling cost and DDIM.}
Equation~\ref{eq:diff-reverse} must be applied $T=1000$ times, and every
application is a full U-Net forward pass, a \emph{network function
evaluation} (NFE). A GAN needs one. Denoising diffusion implicit models
(DDIM) reduce this cost without retraining. They first estimate the clean
image with Equation~\ref{eq:diff-x0-hat} and then re-noise it
deterministically to a lower noise level, reusing the predicted noise
direction:
\begin{equation}
  x_{t'} = \sqrt{\bar\alpha_{t'}}\,\hat x_0
  + \sqrt{1-\bar\alpha_{t'}}\;\varepsilon_\theta(x_t,t),
  \qquad t'<t.
  \label{eq:diff-ddim}
\end{equation}
\noindent
Since no fresh noise enters, the jump from $t$ to $t'$ need not be a single
step. A subsequence of about $50$ time steps gives images of similar quality
to the full $1000$-step chain. The same start noise $x_T$ always yields the
same image, which makes DDIM sampling reproducible.
Table~\ref{tab:diff-train-vs-sample} summarizes training and generation.

\begin{table}[ht]
  \centering
  \begin{tabular}{@{}p{2.5cm}p{5.6cm}p{5.6cm}@{}}
    \toprule
    & Training & Generation \\
    \midrule
    Start & training image $x_0$ & $x_T\sim\mathcal N(0,I)$ \\
    Noise & sampled $\varepsilon$, known & unknown, estimated by
      $\varepsilon_\theta$ \\
    Time steps & one random $t$ per example & all $t=T,\dots,1$ (DDIM: ${\approx}50$) \\
    Network use & one forward and backward pass & one forward pass per step \\
    Formula & Eq.~\ref{eq:diff-marginal}, loss Eq.~\ref{eq:diff-loss} &
      Eq.~\ref{eq:diff-reverse} or Eq.~\ref{eq:diff-ddim} \\
    Result & $\varepsilon_\theta(x_t,t)\approx\varepsilon$ & new image $x_0$ \\
    \bottomrule
  \end{tabular}
  \caption{Training versus generation. Training jumps directly to a random
  noise level; generation must walk the whole chain back.}
  \label{tab:diff-train-vs-sample}
\end{table}

\noindent
In one sentence: the trained object is the network $\varepsilon_\theta$ that
predicts, for every noise level $t$, the noise tensor that should be removed to
move a sample back toward the learned image distribution.

\subsection{Latent Diffusion and Stable Diffusion}

A diffusion model working directly on pixels is expensive. A U-Net must
process a $512\times512\times3$ tensor at every one of many steps, both in
training and in generation. Most of these pixel values encode fine detail that
a good autoencoder can reproduce from a much smaller code. \emph{Latent
diffusion models} (LDM) therefore separate the problem into two stages.

In the first stage, a variational autoencoder with encoder $\mathcal E$ and
decoder $\mathcal D$ is trained to compress and reconstruct images:
$z=\mathcal E(x)\in\mathbb R^{h\times w\times c}$ and
$\tilde x=\mathcal D(z)$. A typical downsampling factor $f=8$ turns
$512\times512\times3$ into $64\times64\times4$. The latent has $64$ times
fewer spatial positions and $48$ times fewer values, yet the reconstruction
is nearly lossless for human perception. In the second stage, the diffusion
process of the previous subsections runs entirely in this latent space. The
loss is Equation~\ref{eq:diff-loss} with $z$ in place of $x$,
\begin{equation}
  \mathcal L_{\mathrm{LDM}} = \mathbb E_{z,\,\varepsilon,\,t}\Bigl[
  \bigl\|\varepsilon-\varepsilon_\theta(z_t,t,\tau_\theta(y))\bigr\|^2\Bigr],
  \qquad z=\mathcal E(x),
  \label{eq:diff-ldm}
\end{equation}
\noindent
where $\tau_\theta(y)$ encodes a condition $y$ such as a text prompt. The
work is thus divided: the VAE performs the actual compression and
decompression, and the U-Net performs multiscale noise prediction on small
tensors. Stable Diffusion is such a latent diffusion model for text-to-image
generation.

\begin{figure}[ht]
  \centering
  \begin{tikzpicture}[
    box/.style={draw, rounded corners=2pt, minimum height=0.85cm,
      align=center, font=\small},
    link/.style={-{Latex[length=2mm]}, thick},
    lbl/.style={font=\footnotesize, align=center},
    region/.style={draw, rounded corners=6pt, thick}
  ]
    \fill[red!6] (-0.9,-2.9) rectangle (1.3,1.35);
    \draw[region] (-0.9,-2.9) rectangle (1.3,1.35);
    \node[lbl, anchor=south] at (0.2,-2.85) {pixel space};
    \fill[green!6] (1.6,-2.9) rectangle (10.3,1.35);
    \draw[region] (1.6,-2.9) rectangle (10.3,1.35);
    \node[lbl, anchor=south] at (6.0,-2.85) {latent space};
    \fill[gray!8] (10.6,-2.9) rectangle (13.9,1.35);
    \draw[region] (10.6,-2.9) rectangle (13.9,1.35);
    \node[lbl, anchor=south] at (12.25,-2.85) {conditioning};

    \node[box] (x) at (0.2,0.6) {$x$};
    \node[box, fill=blue!12] (enc) at (2.5,0.6) {$\mathcal E$};
    \node[box] (z) at (4.3,0.6) {$z$};
    \node[box] (zT) at (8.9,0.6) {$z_T$};
    \draw[link] (x) -- (enc);
    \draw[link] (enc) -- (z);
    \draw[link] (z) -- node[lbl, above] {diffusion process (fixed)} (zT);

    \node[box, fill=orange!25, minimum width=3.2cm] (unet) at (6.6,-1.3)
      {denoising U-Net $\varepsilon_\theta$\\ with cross-attention};
    \node[box] (z0) at (2.5,-1.3) {$z_0$};
    \node[box, fill=blue!12] (dec) at (0.2,-0.6) {$\mathcal D$};
    \node[box] (xt) at (0.2,-1.9) {$\tilde x$};
    \draw[link] (zT.south) -- ++(0,-0.6) -| ($(unet.north)+(1.0,0)$);
    \draw[link] (unet) -- node[lbl, above] {$\times T$ steps} (z0);
    \draw[link] (z0.north) |- (dec.east);
    \draw[link] (dec) -- (xt);

    \node[box] (y) at (12.25,0.6) {prompt $y$};
    \node[box, fill=blue!12] (tau) at (12.25,-1.3) {text encoder\\ $\tau_\theta$};
    \draw[link] (y) -- (tau);
    \draw[link] (tau) -- node[lbl, above] {$K,V$} (unet);
  \end{tikzpicture}
  \caption{Stable Diffusion. The VAE encoder $\mathcal E$ maps an image into
  latent space, where noise is added. The U-Net removes it step by step,
  guided by the text encoding through cross-attention. The decoder
  $\mathcal D$ converts only the final latent $z_0$ back into an image. During
  generation, the upper path is skipped and $z_T$ is drawn from
  $\mathcal N(0,I)$.}
  \label{fig:stable-diffusion}
\end{figure}

\begin{table}[ht]
  \centering
  \begin{tabular}{@{}p{3.1cm}p{4.3cm}p{6.3cm}@{}}
    \toprule
    Component & Input $\to$ output & Responsible for \\
    \midrule
    VAE ($\mathcal E$, $\mathcal D$) & $512^2\times3 \leftrightarrow 64^2\times4$ &
      perceptual quality: edges, textures, fine detail \\
    U-Net $\varepsilon_\theta$ & $(z_t,t,\tau_\theta(y))\to\hat\varepsilon$ &
      semantic content: layout, identity, style \\
    Text encoder $\tau_\theta$ & prompt $\to$ token embeddings & open-vocabulary
      grounding of the prompt \\
    \bottomrule
  \end{tabular}
  \caption{Division of labor in Stable Diffusion. Only the U-Net is iterated;
  the text encoder runs once per prompt and the decoder once per image.}
  \label{tab:sd-components}
\end{table}

\paragraph{Processing flow.}
Figure~\ref{fig:stable-diffusion} shows the three components and how data
flows through them. During training, the VAE encoder compresses each image
into a latent tensor, the forward process noises this latent, and the U-Net
learns to predict the noise, conditioned on the caption of the image. During
generation, the encoder is not needed at all: a random latent $z_T$ is
denoised by the U-Net over many steps, conditioned on the prompt, and the VAE
decoder is used only once, at the very end, to convert the final denoised
latent into an image.

\paragraph{Conditioning by cross-attention.}
The prompt is tokenized and turned into a sequence of embeddings by a
pretrained text encoder $\tau_\theta$; Stable Diffusion uses the CLIP text
encoder, which was trained to align images with their captions. Inside the
U-Net, cross-attention layers (Equation~\ref{eq:transformer-cross-attention})
connect the two modalities. The queries come from the spatial U-Net features,
and keys and values come from the text tokens,
$\operatorname{Attn}(Q,K,V)$ with $Q=W^Q\varphi(z_t)$ and
$K=W^K\tau_\theta(y)$, $V=W^V\tau_\theta(y)$. Every image position can thus
look up the words relevant to it: the region that becomes the dog attends to
``dog,'' the background to ``forest.'' The same mechanism accepts other
conditions, such as semantic maps or reference images, if a suitable encoder
$\tau_\theta$ is used. Table~\ref{tab:sd-components} lists what each
component contributes.

\paragraph{Classifier-free guidance.}
A conditional model $\varepsilon_\theta(z_t,t,y)$ tends to follow the prompt
only loosely. Classifier-free guidance strengthens the condition without any
extra classifier. During training, the condition is replaced by an empty
prompt $\varnothing$ with a probability of about $0.1$, so one network learns
both the conditional and the unconditional prediction. At inference, both are
evaluated and extrapolated:
\begin{equation}
  \tilde\varepsilon = (1+w)\,\varepsilon_\theta(z_t,t,y)
  - w\,\varepsilon_\theta(z_t,t,\varnothing)
  = \varepsilon_\theta(z_t,t,\varnothing)
  + (1+w)\bigl(\varepsilon_\theta(z_t,t,y)-\varepsilon_\theta(z_t,t,\varnothing)\bigr).
  \label{eq:diff-cfg}
\end{equation}
\noindent
The difference between the two predictions is the direction in which the
prompt pushes the image; the guidance scale $w$ amplifies it. With $w=0$ the
plain conditional model is used, with higher diversity but weaker prompt
adherence. Large $w$ makes images very faithful to the prompt but less diverse
and eventually oversaturated or distorted. The default guidance scale of
Stable Diffusion is $7.5$; software libraries usually quote the factor in
front of the difference, i.e.\ $1+w$. Guidance doubles the cost per step,
since two predictions are needed.

\subsection{Fine-Tuning with Low-Rank Adaptation}

A pretrained Stable Diffusion model knows many concepts, but not a particular
dog. Teaching it a new subject from about $100$ photographs by updating all of
its roughly one billion parameters needs much GPU memory and easily destroys
what the model already knows. The transfer strategies of
Section~\ref{sec:transformers} offer a spectrum between full fine-tuning and a
frozen backbone (Table~\ref{tab:adaptation}). \emph{Low-rank adaptation}
(LoRA) keeps the pretrained weights frozen and learns only a small correction
for selected weight matrices, typically the attention projections of the
U-Net.

Let a pretrained layer have weight matrix $W_0\in\mathbb R^{k\times d}$. Full
fine-tuning would update all $k\cdot d$ entries. LoRA instead writes
\begin{equation}
  W = W_0 + \Delta W,\qquad \Delta W = BA,\qquad
  A\in\mathbb R^{r\times d},\quad B\in\mathbb R^{k\times r},\quad
  r\ll\min(d,k),
  \label{eq:lora}
\end{equation}
\noindent
and trains only $A$ and $B$. The layer computes
$Wx=W_0x+B(Ax)$: the input is projected down to $r$ dimensions and back up.
The correction $\Delta W$ has rank at most $r$, so it can change the layer
only along $r$ directions. This restriction works because adapting a
pretrained model to a narrow new task usually needs only a small change of
low intrinsic dimension. $B$ is initialized with zeros, so $\Delta W=0$ at the
start and training begins exactly at the pretrained model. After training,
$BA$ can be added to $W_0$ once, so inference costs no extra time.

\paragraph{Worked example.}
For an attention projection with $d=k=1280$, full fine-tuning updates
$1280^2=1{,}638{,}400$ parameters. LoRA with $r=8$ trains
$r(d+k)=8\cdot2560=20{,}480$ parameters, about $1.25\,\%$ of the full matrix.
The gradients and optimizer state shrink by the same factor, which makes
training feasible on a simple GPU, and the stored result is a file of a few
megabytes that is loaded on top of the unchanged base model.

\begin{table}[ht]
  \centering
  \begin{tabular}{@{}p{3.2cm}p{4.6cm}p{6.0cm}@{}}
    \toprule
    Strategy & What is trained & Typical use \\
    \midrule
    Full fine-tuning & all weights & large target dataset, enough memory \\
    Linear probing & only a new linear head & fixed features, e.g.\
      classification \\
    Prompt/prefix tuning & a few learnable input tokens & steer a frozen
      transformer \\
    LoRA & low-rank $A$, $B$ beside frozen $W_0$ & adapt a large generative
      model on small data and a small GPU \\
    \bottomrule
  \end{tabular}
  \caption{Adaptation strategies for large pretrained models, ordered from
  most to fewest changed weights (prompt tuning and LoRA are both very
  light). LoRA changes the behavior of every adapted layer while training
  only a tiny fraction of the parameters.}
  \label{tab:adaptation}
\end{table}

\paragraph{Practical example: a personalized dog.}
The task is to adapt a pretrained Stable Diffusion v1.5 model so that it
generates new images of one dog, using $100$ photos of it. Every training
image receives the same caption, ``a photo of sks dog.'' The token ``sks'' is
a rare identifier that carries almost no prior meaning, so the model binds the
appearance of this dog to it without overwriting the general concept ``dog.''
Training runs the loss of Equation~\ref{eq:diff-ldm} on the dog photos, but
gradients update only the LoRA matrices. A fine-tuning run looks like this:
\par\noindent\begin{minipage}{\linewidth}
\begin{verbatim}
python finetune_lora_sd.py \
  --pretrained_model stable-diffusion-v1-5/stable-diffusion-v1-5 \
  --train_dir data/dogs --output_dir dog_lora \
  --caption "a photo of sks dog" \
  --validation_prompt "a photo of sks dog running through snow" \
  --max_train_steps 800 --learning_rate 1e-4 --lora_rank 8 \
  --batch_size 1 --gradient_accumulation_steps 4 \
  --mixed_precision
\end{verbatim}
\end{minipage}

\smallskip\noindent
Table~\ref{tab:lora-params} explains the settings. The batch size of $1$ with
$4$ accumulation steps sums the gradients of $4$ images before each update,
so the effective batch size is $4$ while memory holds only one image at a
time. With $800$ update steps the model sees $800\cdot4=3200$ images, i.e.\
about $32$ passes over the $100$ photos. The validation prompt combines the
new token with a scene that never occurs in the training photos; it tests
whether the model has learned the dog as a reusable concept or merely
memorized the training pictures.

\begin{table}[ht]
  \centering
  \begin{tabular}{@{}p{5.8cm}p{8.0cm}@{}}
    \toprule
    Setting & Meaning \\
    \midrule
    \texttt{-{}-caption} & shared caption; binds the new subject to ``sks'' \\
    \texttt{-{}-lora\_rank 8} & rank $r$ of $\Delta W=BA$ \\
    \texttt{-{}-learning\_rate 1e-4} & step size for $A$ and $B$ only \\
    \texttt{-{}-max\_train\_steps 800} & number of parameter updates \\
    \texttt{-{}-batch\_size 1} & images per forward/backward pass \\
    \texttt{-{}-gradient\_accumulation\_steps} & passes summed per update ($4$);
      effective batch size $1\cdot4=4$ \\
    \texttt{-{}-mixed\_precision} & 16-bit arithmetic where safe; less memory,
      faster \\
    \texttt{-{}-seed 123} (prediction) & fixes the start noise $z_T$, so results
      are reproducible \\
    \bottomrule
  \end{tabular}
  \caption{Settings of the LoRA fine-tuning and prediction runs. All of them
  aim at training a large model with the memory of a small GPU.}
  \label{tab:lora-params}
\end{table}

\noindent
For prediction, the LoRA weights are loaded on top of the base model and new
prompts are sampled:
\par\noindent\begin{minipage}{\linewidth}
\begin{verbatim}
python predict_lora_sd.py \
  --pretrained_model stable-diffusion-v1-5/stable-diffusion-v1-5 \
  --lora_dir dog_lora \
  --prompt "a photo of sks dog sitting in a forest, sharp, natural light" \
  --output_dir dog_predictions --num_images 4 --seed 123
\end{verbatim}
\end{minipage}

\smallskip\noindent
Each of the four images starts from a different noise latent, derived from the
fixed seed, and is denoised under the prompt with classifier-free guidance.
The results show dogs in a forest scene that did not appear in the training
set. The pretrained model contributes the forest, the lighting, and the
composition; the LoRA correction contributes the appearance of the dog. If the
generated dogs vary in breed or coloring, the adaptation is still too weak,
for example because of too few steps or too small a rank; if every image
reproduces a training photo, it has overfitted.
